% !TeX root = ../apuntes.tex
\chapter{Tracción y Compresión Uniaxial}

\section{Introducción y Metodología de Análisis}

El estudio de la solicitación axil (tracción y compresión uniaxial) aborda el comportamiento resistente y deformacional de prismas mecánicos rectos sometidos a fuerzas cuya resultante actúa según la directriz longitudinal de la pieza.

El esquema secuencial de resolución de una estructura en Resistencia de Materiales sigue los siguientes pasos:
\begin{enumerate}
    \item \textbf{Modelización de la estructura:} Definición de la geometría mediante prismas mecánicos (barras, tirantes, vigas) y caracterización de las ligaduras y apoyos.
    \item \textbf{Cálculo de solicitaciones exteriores:} Identificación de las cargas aplicadas ($F$) y determinación de las reacciones de enlace mediante las ecuaciones de equilibrio estático de la estática clásica ($\sum \vec{F} = \vec{0}$, $\sum \vec{M} = \vec{0}$).
    \item \textbf{Cálculo de esfuerzos internos seccionales (2D):} Obtención de las leyes de esfuerzo axil ($N_x$), esfuerzo cortante ($Q_y$) y momento flector ($M_{fz}$):
          \begin{equation}
              \{ N_x(x), \quad Q_y(x), \quad M_{fz}(x) \}
          \end{equation}
          En este tema nos centraremos de forma específica en los problemas donde predomina o actúa exclusivamente el \textbf{esfuerzo axil} ($N_x \neq 0$, $Q_y = 0$, $M_{fz} = 0$).
    \item \textbf{Comprobación en las secciones críticas o más desfavorables:}
          \begin{itemize}
              \item \textbf{Criterio de Resistencia (Tensión Máxima):}
                    \begin{equation}
                        \sigma_{\text{máx},d} \le \sigma_{\text{adm,mat}} = f_d = \frac{f_k}{\gamma_m}
                    \end{equation}
              \item \textbf{Criterio de Rigidez / Deformación Excesiva (Deformación Límite):}
                    \begin{equation}
                        \varepsilon_{\text{máx}} \le \varepsilon_{\text{lím}} \quad \Longleftrightarrow \quad \Delta l \le \Delta l_{\text{adm}}
                    \end{equation}
          \end{itemize}
\end{enumerate}

\subsection{Valores Característicos ($k$) frente a Valores de Cálculo ($d$)}

En el cálculo estructural moderno según la normativa técnica (Eurocódigos, CTE, EAE):
\begin{itemize}
    \item Los valores reales esperados o de servicio se denominan \textbf{valores característicos} (subíndice $k$).
    \item Los valores ponderados de seguridad para dimensionamiento se denominan \textbf{valores de cálculo o de diseño} (subíndice $d$).
\end{itemize}

Para las acciones y cargas exteriores se introduce el \textbf{coeficiente de mayoración de cargas} $\gamma_F$ ($\gamma_F \ge 1.0$), de modo que:
\begin{equation}
    F_d = \gamma_F \cdot F_k
\end{equation}

Para las resistencias intrínsecas del material se introduce el \textbf{coeficiente de minoración del material} $\gamma_m$ ($\gamma_m \ge 1.0$), que depende de la naturaleza del material y del control de ejecución:
\begin{itemize}
    \item \textbf{Hormigón:} $\gamma_c \approx 1.50$ ($f_{cd} = f_{ck} / \gamma_c$).
    \item \textbf{Acero estructural:} $\gamma_s \approx 1.05\text{--}1.15$ ($f_{yd} = f_{yk} / \gamma_s$).
    \item \textbf{Madera estructural:} $\gamma_m \approx 1.30$ ($f_{md} = f_{mk} / \gamma_m$).
\end{itemize}

La condición de comprobación resistente se formula como:
\begin{equation}
    \sigma_{\text{máx},d} = \gamma_F \cdot \sigma_k \le \sigma_{\text{adm,mat}} = \frac{f_k}{\gamma_m}
\end{equation}

\begin{ImportanteBox}[title={Advertencia Crítica: Mayoración de Cargas vs. Cálculo de Deformaciones}]
    \textbf{¡OJO! Regla fundamental en dimensionamiento y examen:}
    \begin{itemize}
        \item \textbf{Para las TENSIONES ($\sigma$):} Se emplean los valores de los \textbf{esfuerzos mayorados} ($N_d = \gamma_F N_k$), comprobando frente a la resistencia de cálculo minorada ($\sigma_{\text{máx},d} \le f_d$).
        \item \textbf{Para las DEFORMACIONES y ALARGAMIENTOS ($\varepsilon, \Delta l$):} Se calculan \textbf{SIN MAYORAR}, es decir, en \textbf{Estado Límite de Servicio (ELS)} utilizando directamente los valores característicos de las cargas ($F_k, N_k$).
    \end{itemize}
    \textit{Estrategia recomendada:} Es preferible resolver la estática con las cargas características $F_k$ (obteniendo $N_k, \delta$), calcular los alargamientos con $N_k$, y mayorar al final multiplicando por $\gamma_F$ para verificar la tensión resistente $\sigma_d$.
\end{ImportanteBox}

\begin{figure}[H]
    \centering
    \adjustbox{max width=0.95\textwidth}{%
        \begin{tikzpicture}[>=Latex, node distance=1.8cm, auto]
            % Diagrama de flujo de cálculo con coeficientes de seguridad
            \node [draw, rectangle, rounded corners, fill=blue!10, text width=3.5cm, align=center, thick] (cargas) {\textbf{Cargas}\\\textbf{Características}\\$F_k$};
            \node [draw, rectangle, rounded corners, fill=yellow!15, text width=3.5cm, align=center, thick, right=1.8cm of cargas] (esfuerzos) {\textbf{Esfuerzos Axiles}\\$N_k, Q_k, M_k$};
            \node [draw, rectangle, rounded corners, fill=green!15, text width=3.8cm, align=center, thick, above right=0.8cm and 1.8cm of esfuerzos] (deform) {\textbf{Alargamientos (ELS)}\\$\Delta l = \frac{N_k l}{E \Omega} \le \Delta l_{\text{lím}}$\\\small(Sin mayorar)};
            \node [draw, rectangle, rounded corners, fill=red!10, text width=3.8cm, align=center, thick, below right=0.8cm and 1.8cm of esfuerzos] (tension) {\textbf{Tensiones de Cálculo (ELU)}\\$\sigma_d = \gamma_F \cdot \frac{N_k}{\Omega} \le \frac{f_k}{\gamma_m}$\\\small(Mayoradas con $\gamma_F$)};

            \draw [->, line width=1.3pt] (cargas) -- node[above, font=\footnotesize] {Equilibrio} (esfuerzos);
            \draw [->, line width=1.3pt] (esfuerzos) |- (deform);
            \draw [->, line width=1.3pt] (esfuerzos) |- node[pos=0.7, below, font=\footnotesize] {$\gamma_F$} (tension);
        \end{tikzpicture}%
    }
    \caption{Ruta de cálculo diferenciada para comprobaciones de rigidez (deformaciones sin mayorar) y resistencia (tensiones de cálculo mayoradas).}
    \label{fig:ruta_calculo_seguridad}
\end{figure}

%=============================================================================
\section{Tensiones en Tracción y Compresión Uniaxial}

\subsection{Esfuerzos Seccionales e Integrales de Tensión}

Consideremos un prisma mecánico en equilibrio estático cortado por una sección recta transversal $\Omega$ de normal exterior $\vec{n} = \vec{i}$ orientada según el eje longitudinal $x$. La reducción del sistema de fuerzas interiores sobre el centro de gravedad $G$ define el vector resultante $\vec{R}$ y el momento resultante $\vec{M}_R$:
\begin{align}
    \vec{R}   & = N_x \vec{i} + Q_y \vec{j} + Q_z \vec{k}     \\
    \vec{M}_R & = T \vec{i} + M_{fy} \vec{j} + M_{fz} \vec{k}
\end{align}

Las componentes de reducción se relacionan de forma exacta con la distribución continua de tensiones $\{\sigma_{nx}, \tau_{xy}, \tau_{xz}\}$ mediante las integrales de sección:
\begin{align}
    N_x    & = \iint_{\Omega} \sigma_{nx} \, d\Omega           \\
    Q_y    & = \iint_{\Omega} \tau_{xy} \, d\Omega             \\
    M_{fz} & = \iint_{\Omega} - y \cdot \sigma_{nx} \, d\Omega
\end{align}

En el caso específico de \textbf{solicitación axil pura}:
\begin{equation}
    N_x \neq 0, \qquad Q_y = 0, \qquad M_{fz} = 0, \qquad T = 0
\end{equation}

\begin{figure}[H]
    \centering
    \adjustbox{max width=0.85\textwidth}{%
        \begin{tikzpicture}[>=Latex, scale=1.1]
            % Cilindro / prisma mecánico tridimensional
            \draw[thick, fill=blue!5] (0,-1) arc (-90:90:0.4 and 1) -- (5,1) arc (90:-90:0.4 and 1) -- cycle;
            \draw[thick, fill=blue!15] (5,0) ellipse (0.4 and 1);
            \draw[dashed] (0,-1) arc (-90:90:0.4 and 1);
            \draw[thick] (0,-1) arc (270:90:0.4 and 1);

            % Eje longitudinal x
            \draw[dashdotted, gray!80, line width=1pt] (-0.8,0) -- (6.5,0) node[right, black] {$x$};

            % Sección intermedia de corte
            \draw[thick, fill=orange!25, fill opacity=0.8] (2.8,0) ellipse (0.4 and 1);
            \node[above=1.25cm, font=\small\bfseries, text=orange!80!black] at (2.8,0) {Sección transversal $\Omega$};

            % Distribución uniforme de tensiones sigma_nx
            \foreach \y in {-0.6, -0.3, 0.0, 0.3, 0.6} {
                    \draw[->, line width=1.1pt, red!85!black] (2.8,\y) -- ++(1.2,0);
                }
            \node[right, font=\small\bfseries, text=red!85!black] at (4.0, 0.3) {$\sigma_{nx}$};

            % Resultante N_x (con suficiente separación del eje)
            \draw[->, line width=2pt, blue!90!black] (5.4,0) -- ++(1.4,0) node[right, font=\bfseries] {$N_x$};
            \draw[<-, line width=2pt, blue!90!black] (-0.4,0) -- ++(-1.4,0) node[left, font=\bfseries] {$N_x$};
        \end{tikzpicture}%
    }
    \caption{Prisma mecánico bajo tracción pura y distribución continua de tensiones normales $\sigma_{nx}$ sobre la sección transversal.}
    \label{fig:prisma_axil_3d}
\end{figure}

\subsection{Cálculo de Tensiones y la Hipótesis de Bernoulli-Navier}

Para deducir la distribución analítica de tensiones normales $\sigma_{nx}$ a partir del esfuerzo axil $N_x$, se establecen las hipótesis fundamentales de la Resistencia de Materiales:
\begin{enumerate}
    \item \textbf{Material Elástico y Lineal (Ley de Hooke):}
          \begin{equation}
              \sigma_{nx} = E \cdot \varepsilon_x
          \end{equation}
          donde $E$ es el módulo de elasticidad longitudinal o módulo de Young.
    \item \textbf{Hipótesis de Bernoulli-Navier:}
          Las secciones transversales planas y perpendiculares a la directriz antes de la deformación \textbf{permanecen planas y perpendiculares} a la directriz deformada, sin distorsión apreciable en su propio plano.
\end{enumerate}

Como consecuencia de la hipótesis de Bernoulli, todas las fibras longitudinales experimentan exactamente el mismo alargamiento relativo:
\begin{equation}
    \varepsilon_x = \frac{\Delta l_x}{l_0} = \text{cte} \quad \forall (y,z) \in \Omega
\end{equation}

Sustituyendo en la ley de Hooke, la tensión normal resulta rigurosamente constante en todos los puntos de la sección transversal:
\begin{equation}
    \sigma_{nx} = E \cdot \varepsilon_x = \text{cte}
\end{equation}

Al ser $\sigma_{nx}$ constante en toda el área $\Omega$, puede extraerse fuera del signo integral del esfuerzo axil:
\begin{equation}
    N_x = \iint_{\Omega} \sigma_{nx} \, d\Omega = \sigma_{nx} \iint_{\Omega} d\Omega = \sigma_{nx} \cdot \Omega
\end{equation}

Despejando, se obtiene la expresión fundamental de la \textbf{tensión en tracción y compresión uniaxial}:
\begin{equation}
    \boxed{\sigma_{nx} = \frac{N_x}{\Omega}}
\end{equation}

\begin{figure}[H]
    \centering
    \adjustbox{max width=0.88\textwidth}{%
        \begin{tikzpicture}[>=Latex, scale=1.1]
            % Barra suspendida con peso y deformación
            \draw[line width=1.5pt] (-1, 3.5) -- (1, 3.5);
            \foreach \x in {-0.9, -0.6, -0.3, 0.0, 0.3, 0.6, 0.9} {
                    \draw[thin] (\x, 3.5) -- ++(-0.2, 0.3);
                }

            % Prisma original (línea discontinua)
            \draw[thick, fill=blue!5] (-0.8, 3.5) rectangle (0.8, 0.5);

            % Fibras longitudinales antes y después
            \draw[dashed, red!70!black, thick, fill=red!5] (-0.8, 0.5) rectangle (0.8, -0.3);
            \draw[<->, line width=1pt, red!85!black] (1.1, 0.5) -- node[right, font=\small\bfseries] {Alargamiento $\Delta l$} (1.1, -0.3);
            \draw[<->, line width=1pt] (-1.1, 3.5) -- node[left, font=\small] {$l_0$} (-1.1, 0.5);

            % Carga aplicada
            \draw[->, line width=2pt, blue!90!black] (0, -0.3) -- (0, -1.5) node[below, font=\bfseries] {$F = 10\text{ kN}$};

            % Diagrama de tensiones en la sección (lado derecho)
            \begin{scope}[shift={(4.5, 0.8)}]
                \draw[thick, fill=gray!15] (-0.8, -1.0) rectangle (0.8, 1.0);
                \draw[->, line width=1pt] (-1.5, 0) -- (2.5, 0) node[right] {$x$};
                \draw[->, line width=1pt] (0, -1.4) -- (0, 1.4) node[above] {$y$};

                % Bloque uniforme de tensiones
                \foreach \y in {-0.7, -0.35, 0.0, 0.35, 0.7} {
                        \draw[->, line width=1.2pt, red!85!black] (0.8, \y) -- ++(1.0, 0);
                    }
                \draw[line width=1.2pt, red!85!black] (1.8, -0.7) -- (1.8, 0.7);
                \node[right, font=\small\bfseries, text=red!85!black] at (1.8, -0.35) {$\sigma_{nx} = \frac{N}{\Omega} = \text{cte}$};
            \end{scope}
        \end{tikzpicture}%
    }
    \caption{Tracción simple: conservación de secciones planas (Bernoulli), alargamiento homogéneo de fibras y distribución tensional uniforme.}
    \label{fig:bernoulli_traccion}
\end{figure}

\subsection{Criterio de Signos para Esfuerzo Axil y Tensión Normal}

El signo del esfuerzo axil y de la tensión normal se corresponden unívocamente:
\begin{itemize}
    \item \textbf{Esfuerzo Axil Positivo ($N_x > 0$):} El vector resultante sale de la sección, produciendo \textbf{Tracción} ($\sigma_{nx} > 0$, alargamiento de fibras).
    \item \textbf{Esfuerzo Axil Negativo ($N_x < 0$):} El vector resultante entra hacia la sección, produciendo \textbf{Compresión} ($\sigma_{nx} < 0$, acortamiento de fibras).
\end{itemize}

%=============================================================================
\section{Deformaciones y Desplazamientos en Solicitación Axil}

\subsection{Respuesta Cinemática según el Tipo de Solicitación}

Cada esfuerzo interno induce un mecanismo específico de deformación en el medio continuo:
\begin{table}[H]
    \centering
    \caption{Correspondencia entre esfuerzos seccionales, tensiones y tipos de deformación.}
    \adjustbox{max width=\textwidth}{%
        \renewcommand{\arraystretch}{1.4}
        \begin{tabular}{c c l l}
            \toprule
            \textbf{Esfuerzo Interno}           & \textbf{Tensión Asociada} & \textbf{Tipo de Deformación}  & \textbf{Efecto Físico en las Fibras}                                 \\
            \midrule
            \textbf{Esfuerzo Axil ($N_x$)}      & $\sigma_{nx}$ (Normal)    & Cambio de tamaño longitudinal & Alargamiento / Acortamiento ($\varepsilon_x = \frac{\Delta l}{l_0}$) \\
            \textbf{Esfuerzo Cortante ($Q_y$)}  & $\tau_{xy}$ (Tangencial)  & Cambio de forma / distorsión  & Variación angular ($\gamma_{xy}$) sin cambio de volumen              \\
            \textbf{Momento Flector ($M_{fz}$)} & $\sigma_{nx}$ (Normal)    & Cambio de tamaño variable     & Tracción y compresión simultánea (curvatura seccional)               \\
            \bottomrule
        \end{tabular}%
    }
    \label{tab:esfuerzos_deformaciones}
\end{table}

\begin{figure}[H]
    \centering
    \adjustbox{max width=0.92\textwidth}{%
        \begin{tikzpicture}[>=Latex, scale=1.0]
            % 1. Axil: Alargamiento / Acortamiento
            \node[font=\bfseries] at (1.5, 1.8) {Esfuerzo Axil ($N_x \leftrightarrow \sigma_{nx}$)};
            \draw[thick, fill=blue!10] (0, 0.5) rectangle (3, 1.3);
            \draw[dashed, red, thick] (-0.3, 0.6) rectangle (3.3, 1.2);
            \draw[->, very thick, blue!80!black] (3, 0.9) -- (4, 0.9) node[right] {$N$};
            \draw[<-, very thick, blue!80!black] (-1, 0.9) -- (0, 0.9);
            \node[font=\footnotesize, align=center] at (1.5, 0.1) {Cambio de tamaño: $\varepsilon_x = \frac{\Delta l}{l_0}$};

            % 2. Cortante: Distorsión angular
            \begin{scope}[shift={(6.8, 0)}]
                \node[font=\bfseries] at (1.8, 1.8) {Esfuerzo Cortante ($Q_y \leftrightarrow \tau_{xy}$)};
                \draw[thick, fill=blue!10] (0, 0.5) rectangle (2.5, 1.5);
                \draw[dashed, red, thick] (0.4, 0.5) -- (2.9, 0.5) -- (2.5, 1.5) -- (0, 1.5) -- cycle;
                \draw[->, very thick, orange!90!black] (2.5, 1.5) -- (2.5, 2.3) node[above] {$Q$};
                \node[font=\footnotesize, align=center] at (1.5, 0.1) {Distorsión angular: $\gamma_{xy}$};
            \end{scope}

            % 3. Flector: Curvatura
            \begin{scope}[shift={(13.2, 0)}]
                \node[font=\bfseries] at (1.5, 1.8) {Momento Flector ($M_{fz} \leftrightarrow \sigma_{nx}$)};
                \draw[thick, fill=blue!10] (0, 0.5) rectangle (3, 1.3);
                \draw[dashed, red, thick] (0.2, 1.4) to[bend left=15] (2.8, 1.4) -- (3.0, 0.4) to[bend right=15] (0.0, 0.4) -- cycle;
                \draw[->, very thick, green!60!black] (3.4, 0.9) arc (0:120:0.5) node[right=4pt] {$M_f$};
                \node[font=\footnotesize, align=center] at (1.5, 0.1) {Curvatura: Tracción + Compresión};
            \end{scope}
        \end{tikzpicture}%
    }
    \caption{Comparativa cinemática de las deformaciones producidas por esfuerzo axil, cortante y momento flector.}
    \label{fig:mecanismos_deformacion}
\end{figure}

\subsection{Deducción Diferencial del Alargamiento Total}

Consideremos una barra elástica de longitud $l$ empotrada en su extremo superior $A$ y sometida a una carga axial $F$ en su extremo inferior $B$. Si aislamos un elemento diferencial de longitud inicial $dx$ comprendido entre dos secciones rectas transversales a cota $x$:
\begin{itemize}
    \item La sección a cota $x$ experimenta un desplazamiento longitudinal $u(x)$.
    \item La sección a cota $x + dx$ experimenta un desplazamiento $u(x) + du_x$.
    \item El alargamiento diferencial del elemento $dx$ es $du_x$.
\end{itemize}

Por definición, la deformación longitudinal unitaria $\varepsilon_x$ representa el cociente entre el alargamiento experimentado por el elemento y su longitud inicial:
\begin{equation}
    \varepsilon_x = \frac{du_x}{dx}
\end{equation}

Despejando el diferencial de desplazamiento y aplicando la ley de Hooke uniaxil ($\sigma_{nx} = E \varepsilon_x$):
\begin{equation}
    du_x = \varepsilon_x \, dx = \frac{\sigma_{nx}}{E} \, dx = \frac{N_x(x)}{E(x) \cdot \Omega(x)} \, dx
\end{equation}

El alargamiento o acortamiento total de un tramo continuo comprendido entre las secciones $A$ y $B$ resulta de la suma acumulada (integral) de las deformaciones de todos sus elementos diferenciales:
\begin{equation}
    \boxed{\Delta l_{AB} = \int_{x_A}^{x_B} du_x = \int_{x_A}^{x_B} \varepsilon_x \, dx = \int_{x_A}^{x_B} \frac{N_x(x)}{E(x) \cdot \Omega(x)} \, dx}
\end{equation}

\subsubsection{Caso Fundamental: Barra Homogénea de Sección y Esfuerzo Constantes}
Si a lo largo del tramo $AB$ el esfuerzo axil es constante ($N(x) = N_{AB} = \text{cte}$), la sección transversal no varía ($\Omega(x) = \Omega = \text{cte}$) y el material es uniforme ($E(x) = E = \text{cte}$), todos los términos salen fuera del signo integral:
\begin{equation}
    \Delta l_{AB} = \frac{N_{AB}}{E \cdot \Omega} \int_{x_A}^{x_B} dx = \frac{N_{AB} \cdot (x_B - x_A)}{E \cdot \Omega} = \boxed{\frac{N_{AB} \cdot l_{AB}}{E \cdot \Omega}}
\end{equation}

\begin{figure}[H]
    \centering
    \adjustbox{max width=0.88\textwidth}{%
        \begin{tikzpicture}[>=Latex, scale=1.1]
            % Techo empotrado
            \draw[line width=1.5pt] (-1.2, 4.0) -- (1.2, 4.0);
            \foreach \x in {-1.0, -0.7, -0.4, -0.1, 0.2, 0.5, 0.8, 1.1} {
                    \draw[thin] (\x, 4.0) -- ++(-0.2, 0.3);
                }
            \node[above right, font=\bfseries] at (0, 4.0) {$A$ (Origen $x_A=0$)};

            % Barra prismática vertical
            \draw[thick, fill=blue!8] (-0.6, 4.0) rectangle (0.6, 0.8);
            \node[left, font=\small] at (-0.7, 2.4) {$\Omega = 2\text{ cm}^2$};
            \draw[<->, line width=1pt] (1.0, 4.0) -- node[right, font=\small\bfseries] {$l$} (1.0, 0.8);

            % Elemento diferencial dx
            \draw[thick, fill=red!20] (-0.6, 2.8) rectangle (0.6, 2.4);
            \draw[<->, line width=0.9pt, red!85!black] (-0.9, 2.8) -- node[left, font=\footnotesize\bfseries] {$dx$} (-0.9, 2.4);
            \draw[<->, line width=0.9pt, blue!85!black] (-1.8, 2.4) -- node[left, font=\footnotesize\bfseries] {$du_x$} (-1.8, 2.15);

            % Barra alargada
            \draw[dashed, red!80!black, thick, fill=red!5] (-0.6, 0.8) rectangle (0.6, 0.1);
            \draw[<->, line width=1.1pt, red!85!black] (1.0, 0.8) -- node[right, font=\small\bfseries] {$\Delta l_{AB}$} (1.0, 0.1);
            \node[right, font=\bfseries] at (0.6, 0.1) {$B$ ($x_B=l$)};

            % Carga F en el extremo
            \draw[->, line width=2pt, blue!90!black] (0, 0.1) -- (0, -1.0) node[below, font=\bfseries] {$F = 10\text{ kN}$};

            % Explicación conceptual desplazada a la derecha
            \node[draw, rounded corners, fill=yellow!10, text width=6.2cm, right=3.2cm of {0, 2.2}, font=\small] {
                \textbf{Interpretación Física:}\\
                $\bullet$ Las secciones próximas a $A$ apenas se desplazan.\\
                $\bullet$ Las secciones inferiores acumulan la suma de todas las deformaciones de los elementos superiores.\\
                $\bullet$ El desplazamiento del extremo $B$ es el alargamiento total acumulado:
                \[ u_B = \Delta l_{AB} = \int_{0}^{l} \frac{N}{E\Omega} \, dx \]
            };
        \end{tikzpicture}%
    }
    \caption{Modelización diferencial del alargamiento de una barra prismática bajo carga axil.}
    \label{fig:alargamiento_diferencial}
\end{figure}

\subsection{Efecto Poisson y Deformaciones Transversales}

Cuando una pieza prismática se alarga longitudinalmente en la dirección del axil ($x$), se contrae transversalmente en las direcciones perpendiculares ($y, z$); recíprocamente, cuando se comprime longitudinalmente, se expande transversalmente.

La relación entre las deformaciones transversales unitarias y la deformación longitudinal viene gobernada por el \textbf{coeficiente de Poisson} $\nu$:
\begin{align}
    \varepsilon_x                             & = \frac{\Delta l_x}{l_{0x}} \\
    \varepsilon_y = \frac{\Delta l_y}{l_{0y}} & = - \nu \cdot \varepsilon_x \\
    \varepsilon_z = \frac{\Delta l_z}{l_{0z}} & = - \nu \cdot \varepsilon_x
\end{align}
El signo negativo garantiza la concordancia física: a tracción ($\varepsilon_x > 0$), las secciones transversales adelgazan ($\varepsilon_y, \varepsilon_z < 0$).

\begin{ImportanteBox}[title={¿Por qué se desprecia el efecto Poisson en el cálculo de vigas y barras?}]
    En elementos estructurales prismáticos tipo barra o viga, la longitud $l$ es varios órdenes de magnitud superior al canto o ancho de la sección ($l \gg b, h$). Dado que el coeficiente de Poisson es siempre inferior a la unidad ($\nu < 0.5$, típicamente $\nu \approx 0.20$ en hormigón y $\nu \approx 0.30$ en acero):
    \begin{itemize}
        \item Si una barra se alarga longitudinalmente $\Delta l_x = 1\text{ mm}$, su dimensión transversal apenas varía unas pocas milésimas de milímetro ($\Delta b \approx 0.001\text{ mm}$).
        \item Por consiguiente, en Resistencia de Materiales clásica se asume que \textbf{el área transversal $\Omega$ permanece constante} frente al cálculo de esfuerzos y tensiones sin incurrir en error apreciable.
    \end{itemize}
\end{ImportanteBox}

%=============================================================================
\section{Problema de Aplicación: Dimensionamiento y Comprobación de un Pilar}

\begin{tcolorbox}[colback=blue!5!white,colframe=blue!75!black,title=\textbf{Problema 6 (Ejemplo de Clase): Pilar de Hormigón en Compresión}]
    Se proyecta un pilar vertical de hormigón en masa empotrado en su base y libre en su cabeza, de sección transversal cuadrada de $30\text{ cm} \times 30\text{ cm}$ y longitud $l = 4\text{ m}$. El pilar soporta en su coronación una carga axial característica de servicio $F_k = 600\text{ kN}$ aplicada a compresión.

    \textbf{Datos de los materiales y coeficientes de seguridad:}
    \begin{itemize}
        \item Resistencia característica a compresión del hormigón: $f_{ck} = 25\text{ N/mm}^2 = 25\text{ MPa}$.
        \item Coeficiente de minoración del hormigón: $\gamma_c = 1.50$.
        \item Coeficiente de mayoración de acciones exteriores: $\gamma_F = 1.35$.
        \item Módulo de elasticidad longitudinal del hormigón: $E = 21.5\text{ GPa} = 21500\text{ MPa}$.
    \end{itemize}

    \textbf{Cuestiones a resolver:}
    \begin{enumerate}[label=\textbf{\Alph*)}]
        \item Determinar el valor de cálculo del esfuerzo axil $N_d$ a lo largo del pilar.
        \item Calcular la distribución de tensiones normales en la cabeza del pilar tanto en servicio ($\sigma_k$) como de cálculo ($\sigma_d$).
        \item Comprobar la resistencia del pilar (Estado Límite Último, ELU).
        \item Determinar el acortamiento total $\Delta l$ del pilar bajo la carga de servicio (Estado Límite de Servicio, ELS).
    \end{enumerate}
\end{tcolorbox}

\subsection*{Resolución Detallada}

\subsubsection*{Paso 1: Modelización y Diagrama de Cuerpo Libre (DCL)}
Tumbamos el pilar orientando el eje longitudinal $x$ de $A$ (base empotrada) a $B$ (cabeza libre cargada):
\begin{itemize}
    \item Extremo $A$ ($x = 0$): Empotramiento que restringe desplazamientos y giro. Reacciones: fuerza horizontal $H_A$, fuerza vertical $V_A$ y momento de empotramiento $M_A$.
    \item Extremo $B$ ($x = 4\text{ m}$): Cabeza del pilar sometida a la fuerza axial compresiva $F_k = 600\text{ kN}$ hacia la izquierda ($-\vec{i}$).
\end{itemize}

\begin{figure}[H]
    \centering
    \adjustbox{max width=0.88\textwidth}{%
        \begin{tikzpicture}[>=Latex, scale=1.0]
            % Pilar vertical 3D a la izquierda
            \begin{scope}[shift={(0,0)}]
                \draw[thick, fill=gray!25] (0,0) -- (1.5,0.7) -- (1.5,4.7) -- (0,4) -- cycle;
                \draw[thick, fill=gray!15] (0,4) -- (1.5,4.7) -- (0.8,5.1) -- (-0.7,4.4) -- cycle;
                \draw[thick, fill=gray!35] (0,0) -- (-0.7,0.4) -- (-0.7,4.4) -- (0,4) -- cycle;

                % Base empotrada horizontal con rayado hacia abajo
                \draw[line width=1.5pt] (-1.0, -0.05) -- (1.8, -0.05);
                \foreach \x in {-0.8, -0.5, -0.2, 0.1, 0.4, 0.7, 1.0, 1.3, 1.6} {
                        \draw[thick] (\x, -0.05) -- ++(-0.2, -0.25);
                    }

                \draw[<->, line width=1pt] (-1.1, 0.4) -- node[left] {$l = 4\text{ m}$} (-1.1, 4.4);
                \draw[->, line width=2pt, red!85!black] (0.4, 6.0) -- (0.4, 4.7) node[midway, right=3pt, font=\bfseries] {$F_k = 600\text{ kN}$};
                \node[below=0.35cm, font=\small\bfseries] at (0.4, -0.2) {Sección $30\text{ cm} \times 30\text{ cm}$};
            \end{scope}

            % DCL del pilar horizontal a la derecha
            \begin{scope}[shift={(5.5, 2.5)}]
                \draw[line width=1.5pt] (0, -1.0) -- (0, 1.0);
                \foreach \y in {-0.8, -0.5, -0.2, 0.1, 0.4, 0.7} {
                        \draw[thin] (0, \y) -- ++(-0.3, 0.2);
                    }
                \node[above left] at (0, 0.2) {$A$};

                \draw[very thick, blue!80!black] (0, 0) -- (5, 0) node[above right] {$B$};
                \draw[<->, line width=0.9pt] (0, -0.6) -- node[below] {$4\text{ m}$} (5, -0.6);

                % Reacciones en A: separadas y con el arco de momento en el cuadrante libre superior-derecho
                \draw[<-, line width=1.5pt, blue!90!black] (-1.1, 0) -- (0, 0) node[pos=0.2, below] {$H_A$};
                \draw[->, line width=1.5pt, blue!90!black] (0, -1.1) -- (0, 0) node[midway, left] {$V_A$};
                \draw[->, line width=1.4pt, green!60!black] ($(0,0)+(20:0.65)$) arc (20:120:0.65);
                \node[above right=1pt, green!60!black, font=\bfseries] at ($(0,0)+(60:0.7)$) {$M_A$};

                % Carga en B
                \draw[<-, line width=2pt, red!85!black] (5, 0) -- (6.5, 0) node[right, font=\bfseries] {$F_k = 600\text{ kN}$};
            \end{scope}
        \end{tikzpicture}%
    }
    \caption{Modelización tridimensional y Diagrama de Cuerpo Libre (DCL) del pilar sometido a compresión axial.}
    \label{fig:pilar_dcl}
\end{figure}

\subsubsection*{Paso 2: Ecuaciones de Equilibrio Global}
Aplicando las condiciones de equilibrio estático:
\begin{align}
    \sum F_x = 0 & \implies H_A - F_k = 0 \implies \mathbf{H_A = 600\text{ kN}} \\
    \sum F_y = 0 & \implies \mathbf{V_A = 0}                                    \\
    \sum M_A = 0 & \implies \mathbf{M_A = 0}
\end{align}

\subsubsection*{Paso 3: Esfuerzos Internos por el Método de las Secciones}
Realizando un corte transversal a distancia genérica $x$ ($0 \le x \le 4\text{ m}$) y considerando el equilibrio del tramo izquierdo:
\begin{equation}
    \sum F_x = 0 \implies H_A + N_x(x) = 0 \implies N_x(x) = -H_A = -600\text{ kN}
\end{equation}
El esfuerzo es estrictamente constante y negativo en todo el pilar:
\begin{equation}
    N_k(x) = -600\text{ kN} \quad (\text{Compresión}), \qquad Q_y(x) = 0, \qquad M_{fz}(x) = 0
\end{equation}

\subsubsection*{Apartado A: Esfuerzo Axil de Cálculo ($N_d$)}
Mayorando el esfuerzo característico con el coeficiente de seguridad de acciones $\gamma_F = 1.35$:
\begin{equation}
    \mathbf{N_d} = \gamma_F \cdot N_k = 1.35 \cdot (-600\text{ kN}) = \mathbf{-810\text{ kN}}
\end{equation}

\begin{figure}[H]
    \centering
    \adjustbox{max width=0.88\textwidth}{%
        \begin{tikzpicture}[>=Latex, scale=1.0]
            % Diagrama de Axiles Característico
            \begin{scope}[shift={(0,0)}]
                \draw[thick] (0, 1.2) -- (4, 1.2);
                \draw[thick, fill=blue!10] (0, 0) rectangle (4, 0.8);
                \node at (2, 0.4) {$\boldsymbol{N_k = -600\text{ kN}}$};
                \draw[->, line width=1pt] (0.8, 0.4) -- (1.2, 0.4);
                \draw[<-, line width=1pt] (2.8, 0.4) -- (3.2, 0.4);
                \node[below] at (2, -0.1) {Diagrama de Axiles $N_k$ (Servicio)};
            \end{scope}

            % Flecha de mayoración
            \draw[->, very thick] (4.6, 0.4) -- node[above, font=\footnotesize] {$\times \gamma_F = 1.35$} (6.2, 0.4);

            % Diagrama de Axiles de Cálculo
            \begin{scope}[shift={(6.8,0)}]
                \draw[thick] (0, 1.2) -- (4, 1.2);
                \draw[thick, fill=red!15] (0, 0) rectangle (4, 0.8);
                \node at (2, 0.4) {$\boldsymbol{N_d = -810\text{ kN}}$};
                \draw[->, line width=1pt] (0.8, 0.4) -- (1.2, 0.4);
                \draw[<-, line width=1pt] (2.8, 0.4) -- (3.2, 0.4);
                \node[below] at (2, -0.1) {Diagrama de Axiles $N_d$ (Cálculo)};
            \end{scope}
        \end{tikzpicture}%
    }
    \caption{Leyes de esfuerzo axil característico y de cálculo a lo largo de la directriz del pilar.}
    \label{fig:diagramas_axil_pilar}
\end{figure}

\subsubsection*{Apartado B: Distribución de Tensiones en la Sección}
El área transversal del pilar es:
\begin{equation}
    \Omega = 30\text{ cm} \times 30\text{ cm} = 900\text{ cm}^2 = 900 \times 10^2\text{ mm}^2 = 90000\text{ mm}^2
\end{equation}

\begin{itemize}
    \item \textbf{Tensión característica de servicio:}
          \begin{equation}
              \sigma_k = \frac{N_k}{\Omega} = \frac{-600000\text{ N}}{90000\text{ mm}^2} = -6.67\text{ N/mm}^2 = -6.67\text{ MPa}
          \end{equation}
    \item \textbf{Tensión de cálculo mayorada:}
          \begin{equation}
              \sigma_d = \frac{N_d}{\Omega} = \frac{-810000\text{ N}}{90000\text{ mm}^2} = \mathbf{-9.00\text{ N/mm}^2 = -9.00\text{ MPa}}
          \end{equation}
\end{itemize}
La tensión es una compresión uniforme sobre toda la sección de $9\text{ N/mm}^2$.

\subsubsection*{Apartado C: Comprobación de Resistencia (ELU)}
La resistencia admisible de cálculo del hormigón se obtiene minorando la resistencia característica por su coeficiente parcial de seguridad $\gamma_c = 1.50$:
\begin{equation}
    f_{cd} = \sigma_{\text{adm}} = \frac{f_{ck}}{\gamma_c} = \frac{25\text{ N/mm}^2}{1.50} = \mathbf{16.67\text{ N/mm}^2}
\end{equation}

Evaluando la condición de seguridad en valor absoluto:
\begin{equation}
    |\sigma_d| \le \sigma_{\text{adm}} \iff 9.00\text{ N/mm}^2 \le 16.67\text{ N/mm}^2 \quad \boldsymbol{\checkmark \quad \text{\textbf{CUMPLE}}}
\end{equation}
El pilar trabaja al 54\% de su capacidad resistente máxima; por tanto, \textbf{resiste holgadamente la solicitación}.

\subsubsection*{Apartado D: Acortamiento Total en Servicio (ELS)}
\begin{ImportanteBox}[title={Cálculo de Deformaciones: ¡Sin Mayorar!}]
    El acortamiento elástico se produce bajo las condiciones reales de carga. Por tanto, se opera con el axil característico $N_k = -600\text{ kN}$ y el módulo de elasticidad $E$:
    \begin{itemize}
        \item $N_k = -600\text{ kN} = -600000\text{ N}$.
        \item $l = 4\text{ m} = 4000\text{ mm}$.
        \item $E = 21.5\text{ GPa} = 21500\text{ MPa} = 21500\text{ N/mm}^2$.
        \item $\Omega = 90000\text{ mm}^2$.
    \end{itemize}
\end{ImportanteBox}

Aplicando la expresión del alargamiento elástico:
\begin{align}
    \Delta l_{AB} & = \frac{N_{AB} \cdot l_{AB}}{E \cdot \Omega} = \frac{(-600000\text{ N}) \cdot (4000\text{ mm})}{(21500\text{ N/mm}^2) \cdot (90000\text{ mm}^2)} \notag \\
                  & = \frac{-2.4 \times 10^9}{1.935 \times 10^9} = \mathbf{-1.24\text{ mm}}
\end{align}

El signo negativo indica físicamente que el pilar experimenta un **acortamiento total de $1.24\text{ mm}$**, valor perfectamente admisible para la funcionalidad de la estructura.

%=============================================================================
\section{Problema de Aplicación: Efecto Poisson y Variación Dimensional 3D}

\begin{tcolorbox}[colback=blue!5!white,colframe=blue!75!black,title=\textbf{Problema 7 (Ejemplo de Clase): Barra Prismática con Deformación Transversal de Poisson}]
    Una barra prismática recta de sección rectangular de dimensiones iniciales $a = 5\text{ mm}$ y $b = 4\text{ mm}$ y longitud inicial $l_0 = 50\text{ cm} = 500\text{ mm} = 0.5\text{ m}$ se encuentra rígidamente empotrada en su base $A$ y orientada verticalmente según el eje $x$. En su extremo libre superior $B$ se aplica una carga axial de tracción $F_k = 20\text{ kN} = 20 \times 10^3\text{ N}$.

    \textbf{Propiedades elásticas del material:}
    \begin{itemize}
        \item Módulo de elasticidad longitudinal: $E = 2.1 \times 10^5\text{ N/mm}^2 = 210000\text{ MPa}$ (acero estructural).
        \item Coeficiente de Poisson: $\nu = 0.30$.
    \end{itemize}

    \textbf{Cuestiones a resolver:}
    \begin{enumerate}[label=\textbf{\Alph*)}]
        \item Plantear el Diagrama de Cuerpo Libre (DCL), calcular las reacciones en el empotramiento y trazar las leyes de esfuerzos internos seccionales.
        \item Calcular el alargamiento longitudinal total $\Delta l_x$ y la deformación unitaria longitudinal $\varepsilon_x$ según el eje de la pieza.
        \item Determinar las deformaciones transversales unitarias $\varepsilon_y$ y $\varepsilon_z$ debidas al efecto Poisson.
        \item Calcular las dimensiones finales efectivas de la pieza deformada: longitud final $l'$, ancho final $a'$ y espesor final $b'$.
    \end{enumerate}
\end{tcolorbox}

\subsection*{Resolución Detallada}

\subsubsection*{Paso 1: Modelización 3D y Diagrama de Cuerpo Libre (DCL)}
Disponemos el eje longitudinal $x$ a lo largo de la directriz de la barra, desde el empotramiento $A$ ($x = 0$) hasta el extremo libre superior $B$ ($x = 0.5\text{ m}$), con los ejes transversales $y$ y $z$ contenidos en el plano de la sección recta ($5\text{ mm} \times 4\text{ mm}$).

\begin{figure}[H]
    \centering
    \adjustbox{max width=0.98\textwidth}{%
        \begin{tikzpicture}[>=Latex, scale=1.05]
            % -------------------------------------------------------------
            % 1. BARRA VERTICAL 3D (Izquierda)
            % -------------------------------------------------------------
            \begin{scope}[shift={(0,0)}]
                \draw[thick, fill=blue!10] (0,0) -- (0.8,0.4) -- (0.8,4.5) -- (0,4.1) -- cycle;
                \draw[thick, fill=blue!20] (0,4.1) -- (0.8,4.5) -- (0.4,4.7) -- (-0.4,4.3) -- cycle;
                \draw[thick, fill=blue!15] (0,0) -- (-0.4,0.2) -- (-0.4,4.3) -- (0,4.1) -- cycle;

                % Base empotrada horizontal con rayado hacia abajo
                \draw[line width=1.4pt] (-0.7, -0.05) -- (1.1, -0.05);
                \foreach \x in {-0.5, -0.2, 0.1, 0.4, 0.7, 1.0} {
                    \draw[thick] (\x, -0.05) -- ++(-0.2, -0.25);
                }

                \draw[<->, line width=0.9pt] (-0.7, 0.2) -- node[left] {$l_0 = 50\text{ cm}$} (-0.7, 4.3);
                \draw[->, line width=2pt, red!85!black] (0.2, 4.4) -- (0.2, 5.7) node[above, font=\bfseries] {$F_k = 20\text{ kN}$};

                % Ejes locales 3D en la cabeza
                \draw[->, line width=1pt] (0.2, 4.4) -- ++(0, 1.0) node[right, font=\footnotesize] {$x$};
                \draw[->, line width=1pt] (0.2, 4.4) -- ++(0.7, 0.35) node[right, font=\footnotesize] {$y$};
                \draw[->, line width=1pt] (0.2, 4.4) -- ++(0.9, -0.2) node[below right, font=\footnotesize] {$z$};

                \node[below=0.25cm, font=\footnotesize] at (0.2, -0.1) {$a = 5\text{ mm}$, $b = 4\text{ mm}$};
            \end{scope}

            % -------------------------------------------------------------
            % 2. DCL Y DIAGRAMA DE ESFUERZOS AXIL N(x) (Derecha)
            % -------------------------------------------------------------
            \begin{scope}[shift={(5.2, 0)}]
                % --- Sub-bloque superior: DCL ---
                \begin{scope}[shift={(0, 3.5)}]
                    \node[above, font=\bfseries] at (2.5, 1.1) {Diagrama de Cuerpo Libre (DCL)};

                    \draw[line width=1.5pt] (0, -0.9) -- (0, 0.9);
                    \foreach \y in {-0.7, -0.4, -0.1, 0.2, 0.5, 0.8} {
                        \draw[thin] (0, \y) -- ++(-0.3, 0.2);
                    }
                    \node[above left] at (0, 0.2) {$A$};

                    \draw[very thick, blue!80!black] (0, 0) -- (5, 0) node[above right] {$B$};
                    \draw[<->, line width=0.9pt] (0, -0.6) -- node[below] {$l_0 = 0.5\text{ m}$} (5, -0.6);

                    % Reacciones en A
                    \draw[<-, line width=1.5pt, blue!90!black] (-1.1, 0) -- (0, 0) node[pos=0.2, below] {$H_A$};
                    \draw[->, line width=1.5pt, blue!90!black] (0, -1.0) -- (0, 0) node[midway, left] {$V_A$};
                    \draw[->, line width=1.4pt, green!60!black] ($(0,0)+(20:0.6)$) arc (20:120:0.6);
                    \node[above right=1pt, green!60!black, font=\bfseries] at ($(0,0)+(60:0.65)$) {$M_A$};

                    % Carga en B
                    \draw[->, line width=2pt, red!85!black] (5, 0) -- (6.5, 0) node[right, font=\bfseries] {$F_k = 20\text{ kN}$};
                \end{scope}

                % Líneas de proyección vertical punteadas entre DCL y Diagrama N
                \draw[dashed, gray!60] (0, 3.5) -- (0, 0.2);
                \draw[dashed, gray!60] (5, 3.5) -- (5, 0.2);

                % --- Sub-bloque inferior: Diagrama de Axiles N(x) ---
                \begin{scope}[shift={(0, 0.6)}]
                    \node[above, font=\bfseries] at (2.5, 1.2) {Diagrama de Esfuerzos Axiles $N(x)$};

                    % Eje de referencia
                    \draw[thick] (0, 0) -- (5, 0);

                    % Bloque de axil positivo (+20 kN constante)
                    \draw[thick, fill=green!15] (0, 0) rectangle (5, 0.9);
                    \node[font=\bfseries] at (2.5, 0.45) {$N(x) = +20\text{ kN}\quad (\text{Tracción})$};

                    % Flechas de tracción salientes
                    \draw[<-, line width=1pt] (0.6, 0.45) -- (1.0, 0.45);
                    \draw[->, line width=1pt] (4.0, 0.45) -- (4.4, 0.45);

                    \node[below left] at (0, 0) {$A$};
                    \node[below right] at (5, 0) {$B$};
                \end{scope}
            \end{scope}
        \end{tikzpicture}%
    }
    \caption{Modelización tridimensional de la barra prismática junto con su Diagrama de Cuerpo Libre (DCL) y su correspondiente ley de esfuerzos axiles $N(x)$.}
    \label{fig:pilar_poisson_dcl}
\end{figure}

\subsubsection*{Paso 2: Equilibrio Estático y Reacciones}
Tumbando la barra horizontalmente con origen en el empotramiento $A$ ($x = 0$):
\begin{align}
    \sum F_x = 0 & \implies -H_A + F_k = 0 \implies \mathbf{H_A = 20 \times 10^3\text{ N} = 20\text{ kN}} \\
    \sum F_y = 0 & \implies \mathbf{V_A = 0}                                                              \\
    \sum M_A = 0 & \implies \mathbf{M_A = 0}
\end{align}

\subsubsection*{Paso 3: Esfuerzos Internos por el Método de las Secciones}
Realizando un corte a distancia genérica $x$ ($0 \le x \le 0.5\text{ m}$):
\begin{equation}
    \sum F_x = 0 \implies -H_A + N_x(x) = 0 \implies \mathbf{N_x(x) = H_A = +20\text{ kN}} \quad (\text{Tracción})
\end{equation}
El esfuerzo es estrictamente constante y positivo en toda la longitud de la barra, con $Q_y(x) = 0$ y $M_{fz}(x) = 0$.

\subsubsection*{Apartado B: Alargamiento Longitudinal ($\Delta l_x$) y Deformación $\varepsilon_x$}
El área transversal de la sección recta es:
\begin{equation}
    \Omega = a \cdot b = 5\text{ mm} \times 4\text{ mm} = \mathbf{20\text{ mm}^2}
\end{equation}

El alargamiento elástico longitudinal en la dirección $x$ resulta:
\begin{align}
    \Delta l_x & = \frac{N \cdot l_0}{E \cdot \Omega} = \frac{(20 \times 10^3\text{ N}) \cdot (500\text{ mm})}{(2.1 \times 10^5\text{ N/mm}^2) \cdot (20\text{ mm}^2)} \notag \\
               & = \frac{10 \times 10^6\text{ N}\cdot\text{mm}}{4.2 \times 10^6\text{ N}} = \mathbf{2.381\text{ mm} \approx 2.38\text{ mm}}
\end{align}
El desplazamiento del extremo libre es positivo (alargamiento): $\mathbf{\Delta l_x = +2.38\text{ mm} = +0.00238\text{ m}}$.

La deformación unitaria longitudinal $\varepsilon_x$ se calcula como:
\begin{equation}
    \boldsymbol{\varepsilon_x} = \frac{\Delta l_x}{l_{0x}} = \frac{2.381\text{ mm}}{500\text{ mm}} = \mathbf{0.00476} \quad (4.76 \times 10^{-3} = 0.476\text{\%})
\end{equation}

\subsubsection*{Apartado C: Deformaciones Transversales Unitarias ($\varepsilon_y, \varepsilon_z$)}
Aplicando la formulación del efecto Poisson con $\nu = 0.30$:
\begin{equation}
    \boldsymbol{\varepsilon_y = \varepsilon_z} = -\nu \cdot \varepsilon_x = -0.30 \cdot (0.00476) = \mathbf{-0.001428} \quad (-1.428 \times 10^{-3})
\end{equation}
El signo negativo indica que las fibras transversales experimentan un \textbf{acortamiento relativo (contracción lateral)}.

\subsubsection*{Apartado D: Dimensiones Finales de la Pieza Deformada}
A partir de la definición de deformación unitaria en cada eje coordenado ($\varepsilon_i = \frac{\Delta l_i}{l_{0i}}$):

\begin{enumerate}
    \item \textbf{Longitud final ($l'$ según el eje $x$):}
          \begin{equation}
              l' = l_0 + \Delta l_x = 500\text{ mm} + 2.381\text{ mm} = \mathbf{502.381\text{ mm} \approx 502.38\text{ mm}}
          \end{equation}
    \item \textbf{Ancho final ($a'$ según el eje $y$, con $a = 5\text{ mm}$):}
          \begin{align}
              \Delta a & = \varepsilon_y \cdot a = (-0.001428) \cdot (5\text{ mm}) = -0.00714\text{ mm}                        \\
              a'       & = a + \Delta a = 5\text{ mm} - 0.00714\text{ mm} = \mathbf{4.99286\text{ mm} \approx 4.993\text{ mm}}
          \end{align}
    \item \textbf{Espesor final ($b'$ según el eje $z$, con $b = 4\text{ mm}$):}
          \begin{align}
              \Delta b & = \varepsilon_z \cdot b = (-0.001428) \cdot (4\text{ mm}) = -0.00571\text{ mm}                        \\
              b'       & = b + \Delta b = 4\text{ mm} - 0.00571\text{ mm} = \mathbf{3.99429\text{ mm} \approx 3.994\text{ mm}}
          \end{align}
\end{enumerate}

\begin{table}[H]
    \centering
    \caption{Balance comparativo de dimensiones iniciales frente a finales en la barra traccionada.}
    \adjustbox{max width=0.85\textwidth}{%
        \renewcommand{\arraystretch}{1.35}
        \begin{tabular}{l c c c c}
            \toprule
            \textbf{Dimensión} & \textbf{Eje} & \textbf{Valor Inicial ($l_0$)} & \textbf{Variación ($\Delta l$)} & \textbf{Valor Final Deformado ($l'$)}            \\
            \midrule
            Longitud           & $x$          & $500.000\text{ mm}$            & $+2.381\text{ mm}$              & $\mathbf{502.381\text{ mm}}$ (Alargamiento)      \\
            Ancho              & $y$          & $5.000\text{ mm}$              & $-0.007\text{ mm}$              & $\mathbf{4.993\text{ mm}}$ (Contracción lateral) \\
            Espesor            & $z$          & $4.000\text{ mm}$              & $-0.006\text{ mm}$              & $\mathbf{3.994\text{ mm}}$ (Contracción lateral) \\
            \bottomrule
        \end{tabular}%
    }
    \label{tab:dimensiones_finales_poisson}
\end{table}

Este resultado numérico confirma de forma experimental lo expuesto en la teoría: mientras que la longitud experimenta una variación apreciable de $+2.38\text{ mm}$, las dimensiones transversales solo se reducen unas milésimas de milímetro ($\approx -0.007\text{ mm}$), por lo que el área transversal se puede considerar inalterada a efectos de tensiones en piezas esbeltas.

%=============================================================================
\section{Problema de Aplicación: Pieza Compuesta con Sección Escalonada y Variable}

\begin{tcolorbox}[colback=blue!5!white,colframe=blue!75!black,title=\textbf{Problema 8: Barra Compuesta (Tubo de Acero y Tronco de Pirámide) bajo Compresión}]
    Se considera un elemento estructural vertical compuesto por dos tramos diferenciados empotrado en su base $D$:
    \begin{itemize}
        \item \textbf{Tramo superior $ABC$ (Tubo circular hueco de acero S275):} Longitud total $l_{AC} = 1.0\text{ m}$ dividida en dos subtramos $AB$ y $BC$ de $0.5\text{ m}$ cada uno.
              \begin{itemize}
                  \item Diámetro exterior: $\phi_{\text{ext}} = 200\text{ mm}$ ($R_{\text{ext}} = 100\text{ mm}$).
                  \item Espesor de pared: $e = 6\text{ mm}$ ($R_{\text{int}} = 94\text{ mm}$).
                  \item Límite elástico característico del acero: $f_{yk} = 275\text{ N/mm}^2$.
                  \item Coeficiente de seguridad del material: $\gamma_s = 1.05$.
              \end{itemize}
        \item \textbf{Tramo inferior $CD$ (Prisma en tronco de pirámide de fábrica / hormigón):} Longitud $l_{CD} = 1.0\text{ m}$.
              \begin{itemize}
                  \item Sección superior $C$: cuadrada de $250\text{ mm} \times 250\text{ mm}$.
                  \item Sección inferior $D$ (base empotrada): rectangular de $250\text{ mm} \times 300\text{ mm}$.
                  \item Resistencia característica a compresión: $f_{ck} = 20\text{ N/mm}^2$.
                  \item Coeficiente de seguridad del material: $\gamma_c = 1.50$.
              \end{itemize}
    \end{itemize}

    \textbf{Acciones exteriores aplicadas (valores de cálculo ponderados):}
    \begin{itemize}
        \item En el extremo superior $A$: carga puntual axil de compresión $F_{d1} = 120\text{ kN}$ hacia abajo.
        \item En la sección intermedia $B$: carga puntual axil $F_{d2} = 30\text{ kN}$ orientada hacia arriba.
    \end{itemize}

    \textbf{Cuestiones a resolver:}
    \begin{enumerate}[label=\textbf{\Alph*)}]
        \item Determinar las reacciones en la base empotrada $D$ y trazar la ley de esfuerzos axiles de cálculo $N_d(x)$ a lo largo de toda la pieza.
        \item Calcular el área de las secciones transversales clave y determinar la distribución de tensiones normales de cálculo $\sigma_d$ a lo largo de la pieza, representando su diagrama.
        \item Realizar la comprobación de resistencia frente al Estado Límite Último (ELU) en las secciones más desfavorables de cada material.
    \end{enumerate}
\end{tcolorbox}

\subsection*{Resolución Detallada}

\subsection*{a) Ley de Esfuerzos Axiles $N_d$}

\subsubsection*{1. Diagrama de Cuerpo Libre y Cálculo de Reacciones}
Disponemos la pieza horizontalmente a lo largo del eje $x$ con origen en la base empotrada $D$ ($x = 0$) hasta la coronación libre $A$ ($x = 2.0\text{ m}$):
\begin{itemize}
    \item $x = 0$: Empotramiento $D$ con reacciones $H_D$, $V_D$, $M_D$.
    \item $x = 1.0\text{ m}$: Sección $C$ (unión entre el tronco de pirámide y el tubo de acero).
    \item $x = 1.5\text{ m}$: Sección $B$ con carga axial puntual de $30\text{ kN}$ hacia la derecha ($+x$).
    \item $x = 2.0\text{ m}$: Sección $A$ con carga axial puntual de $120\text{ kN}$ hacia la izquierda (compresiva, $-x$).
\end{itemize}

\begin{figure}[H]
    \centering
    \adjustbox{max width=0.88\textwidth}{%
        \begin{tikzpicture}[>=Latex, scale=1.05]
            % Base D en x = 0
            \draw[line width=1.5pt] (0, -1.4) -- (0, 1.4);
            \foreach \y in {-1.2, -0.8, -0.4, 0.0, 0.4, 0.8, 1.2} {
                \draw[thin] (0, \y) -- ++(-0.3, 0.2);
            }
            \node[above left] at (0, 0.1) {$D$};

            % Tronco de pirámide CD (de x=0 a x=3) con corte en zigzag al final
            \draw[thick, fill=gray!15] (0, 1.1) -- (3, 0.6) -- (3, -0.6) -- (0, -1.1) -- cycle;
            
            % Ondas de corte de sección
            \draw[thick, gray!70] (3, -0.7) to[bend left=20] (3, 0) to[bend right=20] (3, 0.7);
            \draw[thick, gray!70] (6, -0.7) to[bend left=20] (6, 0) to[bend right=20] (6, 0.7);

            % Tubo de acero ABC (de x=3 a x=9)
            \draw[thick, fill=blue!10] (3, 0.6) -- (9, 0.6) -- (9, -0.6) -- (3, -0.6) -- cycle;

            % Eje x
            \draw[->, dashdotted, line width=0.8pt] (0,0) -- (10.2, 0) node[right] {$x$};

            % Puntos notables
            \filldraw[red] (3, 0) circle (2pt) node[below=0.6cm, black] {$C$};
            \filldraw[red] (6, 0) circle (2pt) node[below=0.6cm, black] {$B$};
            \filldraw[red] (9, 0) circle (2pt) node[below=0.6cm, black] {$A$};

            % Cotas
            \draw[<->, line width=0.8pt] (0, 1.6) -- node[above] {$1\text{ m}$} (3, 1.6);
            \draw[<->, line width=0.8pt] (3, 1.6) -- node[above] {$0.5$} (6, 1.6);
            \draw[<->, line width=0.8pt] (6, 1.6) -- node[above] {$0.5$} (9, 1.6);

            % Cargas en B y A
            \draw[->, line width=2pt, red!90!black] (6, 0) -- (7.5, 0) node[above, font=\bfseries] {$30$};
            \draw[<-, line width=2.2pt, red!90!black] (9, 0) -- (10.8, 0) node[above, font=\bfseries] {$120$};

            % Reacciones en D
            \draw[->, line width=1.8pt, blue!90!black] (-1.3, 0) -- (0, 0) node[midway, above, font=\bfseries] {$H_D$};
            \draw[->, line width=1.8pt, blue!90!black] (0, -1.3) -- (0, 0) node[midway, left, font=\bfseries] {$V_D$};
            \draw[->, line width=1.5pt, green!60!black] ($(0,0)+(20:0.7)$) arc (20:120:0.7);
            \node[above right=1pt, green!60!black, font=\bfseries] at ($(0,0)+(60:0.75)$) {$M_D$};
        \end{tikzpicture}%
    }
    \caption{Diagrama de Cuerpo Libre (DCL) horizontal y planteamiento de equilibrio global.}
    \label{fig:dcl_horizontal_p8}
\end{figure}

Aplicando las ecuaciones de la estática clásica:
\begin{equation}
    \begin{cases}
        \sum F_H = 0 \implies H_D + 30 - 120 = 0 \implies \mathbf{H_D = 90\text{ kN}} \\
        \sum F_V = 0 \implies \mathbf{V_D = 0} \\
        \sum M_D = 0 \implies \mathbf{M_D = 0}
    \end{cases}
\end{equation}

\subsubsection*{2. Método de las Secciones}
\begin{itemize}
    \item \textbf{Sección 1 ($S.1$): $x \in [0, 1.5\text{ m}]$} (comprende el tramo de fábrica $CD$ y el tramo de acero $BC$):
    \begin{equation}
        \sum F_H = 0 \implies 90 + N = 0 \implies \mathbf{N = -90\text{ kN}} \quad (\text{Compresión})
    \end{equation}

    \item \textbf{Sección 2 ($S.2$): $x \in [1.5, 2\text{ m}]$} (tramo de acero $AB$):
    \begin{equation}
        \sum F_H = 0 \implies 90 + 30 + N = 0 \implies \mathbf{N = -120\text{ kN}} \quad (\text{Compresión})
    \end{equation}
\end{itemize}

\begin{figure}[H]
    \centering
    \adjustbox{max width=0.85\textwidth}{%
        \begin{tikzpicture}[>=Latex, scale=1.0]
            % Eje de referencia horizontal
            \draw[thick] (0, 1.2) -- (9, 1.2);
            \node[above] at (0, 1.2) {$D$};
            \node[above] at (4.5, 1.2) {$C$};
            \node[above] at (6.75, 1.2) {$B$};
            \node[above] at (9, 1.2) {$A$};

            % Diagrama de axiles escalonado hacia abajo (negativo)
            \draw[thick, fill=red!10] (0, 1.2) -- (0, 0.2) -- (6.75, 0.2) -- (6.75, -0.6) -- (9, -0.6) -- (9, 1.2) -- cycle;

            % Valores numéricos
            \node at (3.3, 0.7) {$\boldsymbol{-90}$};
            \draw[->, line width=0.8pt] (4.5, 0.7) -- (5.0, 0.7);
            \draw[<-, line width=0.8pt] (5.5, 0.7) -- (6.0, 0.7);

            \node at (7.87, -0.2) {$\boldsymbol{-120}$};
            \draw[->, line width=0.8pt] (7.2, 0.7) -- (7.7, 0.7);
            \draw[<-, line width=0.8pt] (8.2, 0.7) -- (8.7, 0.7);

            % Discontinuidad en B
            \draw[dashed] (6.75, 1.2) -- (6.75, -0.6);

            \node[left] at (-0.2, 0.2) {$N_d$};
        \end{tikzpicture}%
    }
    \caption{Diagrama de esfuerzos axiles $N_d$ [$\text{kN}$] horizontal.}
    \label{fig:diagrama_axil_horizontal_p8}
\end{figure}

\subsection*{b) Distribución de Tensiones}

\subsubsection*{1. Características Geométricas de las Secciones Transversales}
\begin{itemize}
    \item \textbf{Sección tubo de acero ($ABC$):} Sección anular hueca constante:
    \begin{equation}
        \Omega_A = \Omega_B = \Omega_C = \pi \left( R_{\text{ext}}^2 - R_{\text{int}}^2 \right) = \pi \left( 100^2 - 94^2 \right) = \mathbf{3656.8\text{ mm}^2}
    \end{equation}

    \item \textbf{Sección tronco de pirámide ($CD$):}
    \begin{align}
        \text{En } C: \quad \Omega_C &= 250\text{ mm} \times 250\text{ mm} = \mathbf{62500\text{ mm}^2} \\
        \text{En } D: \quad \Omega_D &= 250\text{ mm} \times 300\text{ mm} = \mathbf{75000\text{ mm}^2}
    \end{align}
\end{itemize}

\begin{figure}[H]
    \centering
    \adjustbox{max width=0.75\textwidth}{%
        \begin{tikzpicture}[>=Latex, scale=1.0]
            % Sección C (250 x 250)
            \begin{scope}[shift={(0,0)}]
                \draw[thick, fill=gray!15] (-1.25, -1.25) rectangle (1.25, 1.25);
                \draw[->, line width=0.8pt] (0, 0) -- (0, 1.8) node[above] {$z$};
                \draw[->, line width=0.8pt] (0, 0) -- (1.8, 0) node[right] {$y$};
                \node at (0, 0) {$C$};
                \draw[<->, line width=0.6pt] (-1.25, -1.5) -- node[below] {$250$} (1.25, -1.5);
                \draw[<->, line width=0.6pt] (-1.5, -1.25) -- node[left] {$250$} (-1.5, 1.25);
                \node[below=0.6cm, font=\small\bfseries] at (0, -1.5) {Sección $C$};
            \end{scope}

            % Sección D (250 x 300)
            \begin{scope}[shift={(5.5,0)}]
                \draw[thick, fill=gray!15] (-1.5, -1.25) rectangle (1.5, 1.25);
                \draw[->, line width=0.8pt] (0, 0) -- (0, 1.8) node[above] {$z$};
                \draw[->, line width=0.8pt] (0, 0) -- (2.1, 0) node[right] {$y$};
                \node at (0, 0) {$D$};
                \draw[<->, line width=0.6pt] (-1.5, -1.5) -- node[below] {$300$} (1.5, -1.5);
                \draw[<->, line width=0.6pt] (-1.8, -1.25) -- node[left] {$250$} (-1.8, 1.25);
                \node[below=0.6cm, font=\small\bfseries] at (0, -1.5) {Sección $D$};
            \end{scope}
        \end{tikzpicture}%
    }
    \caption{Definición de las secciones transversales en los extremos $C$ y $D$ de la obra de fábrica.}
    \label{fig:secciones_transversales_p8}
\end{figure}

\subsubsection*{2. Cálculo de Tensiones Normales}
\begin{align}
    \left.
    \begin{aligned}
        \sigma_{A,d} &= \frac{-120 \cdot 10^3\text{ N}}{3656.8\text{ mm}^2} = \mathbf{-32.81\text{ N/mm}^2} \\[6pt]
        \sigma_{B,d} &= \frac{-90 \cdot 10^3\text{ N}}{3656.8\text{ mm}^2} = \mathbf{-24.61\text{ N/mm}^2}
    \end{aligned}
    \right\} \quad \text{\textbf{acero}} \\[10pt]
    \left.
    \begin{aligned}
        \sigma_{C,d} &= \frac{-90 \cdot 10^3\text{ N}}{62500\text{ mm}^2} = \mathbf{-1.44\text{ N/mm}^2} \\[6pt]
        \sigma_{D,d} &= \frac{-90 \cdot 10^3\text{ N}}{75000\text{ mm}^2} = \mathbf{-1.20\text{ N/mm}^2}
    \end{aligned}
    \right\} \quad \text{\textbf{obra de fábrica}}
\end{align}

\begin{figure}[H]
    \centering
    \adjustbox{max width=0.88\textwidth}{%
        \begin{tikzpicture}[>=Latex, scale=1.0]
            % Eje de referencia horizontal
            \draw[thick] (0, 1.5) -- (9, 1.5);
            \node[above] at (0, 1.5) {$D$};
            \node[above] at (4.5, 1.5) {$C$};
            \node[above] at (6.75, 1.5) {$B$};
            \node[above] at (9, 1.5) {$A$};

            % Trapecio suave en fábrica: de -1.2 a -1.44
            \draw[thick, fill=blue!10] (0, 1.5) -- (0, 1.1) -- (4.5, 0.95) -- (4.5, 1.5) -- cycle;
            \node[left, font=\footnotesize] at (0, 1.1) {$-1.2$};
            \node[above right, font=\footnotesize] at (4.5, 0.95) {$-1.44$};

            % Escalonado en acero: -24.61 y salto en B a -32.81
            \draw[thick, fill=blue!20] (4.5, 1.5) -- (4.5, 0.3) -- (6.75, 0.3) -- (6.75, -0.6) -- (9, -0.6) -- (9, 1.5) -- cycle;
            
            % Discontinuidades y valores
            \draw[dashed] (4.5, 1.5) -- (4.5, 0.3);
            \draw[dashed] (6.75, 1.5) -- (6.75, -0.6);

            \node[below, font=\footnotesize] at (5.6, 0.3) {$-24.61$};
            \draw[->, line width=0.8pt] (5.4, 0.8) -- (5.9, 0.8);
            \draw[<-, line width=0.8pt] (6.2, 0.8) -- (6.7, 0.8);

            \node[below, font=\footnotesize\bfseries] at (7.87, -0.6) {$-32.81$};
            \node[left] at (-0.2, 1.1) {$\sigma_n$};
        \end{tikzpicture}%
    }
    \caption{Diagrama de tensiones normales $\sigma_n$ [$\text{N/mm}^2$] horizontal.}
    \label{fig:diagrama_tensiones_horizontal_p8}
\end{figure}

\subsection*{c) Comprobar la Resistencia de la Pieza}
Verificamos las secciones más desfavorables de cada material frente a sus tensiones admisibles de cálculo:
\begin{enumerate}
    \item \textbf{Tubo de acero $ABC$:}
    \begin{equation}
        |\sigma_{\text{máx,acero}}| = |-32.81\text{ N/mm}^2| \stackrel{?}{\le} f_{yd} = \frac{275}{1.05} = 261.90\text{ N/mm}^2 \quad (\text{o } \le 275) \quad \boldsymbol{\checkmark \quad \text{\textbf{Sí resiste}}}
    \end{equation}

    \item \textbf{Tronco de pirámide de fábrica $CD$:}
    \begin{equation}
        |\sigma_{\text{máx,fáb}}| = |-1.44\text{ N/mm}^2| \stackrel{?}{\le} f_{cd} = \frac{20}{1.50} = 13.33\text{ N/mm}^2 \quad (\text{o } < 20) \quad \boldsymbol{\checkmark \quad \text{\textbf{Sí resiste}}}
    \end{equation}
\end{enumerate}

%=============================================================================
\section{Problema de Aplicación: Influencia del Peso Propio en Solicitación Axil}

\begin{tcolorbox}[colback=blue!5!white,colframe=blue!75!black,title=\textbf{Problema 9: Columna Compuesta Escalonada con Carga Puntual y Peso Propio}]
    Se proyecta un soporte vertical compuesto por dos elementos coaxiales continuos de acero de densidad $\gamma = 7850\text{ kg/m}^3$, empotrado rígidamente en su base inferior:
    \begin{itemize}
        \item \textbf{Tramo superior (Tubo hueco de acero):}
              \begin{itemize}
                  \item Altura: $l_1 = 1.0\text{ m}$.
                  \item Diámetro exterior: $\phi_{\text{ext}} = 150\text{ mm}$ ($R_{\text{ext}} = 75\text{ mm}$).
                  \item Espesor de paredes: $e = 5\text{ mm}$ ($R_{\text{int}} = 70\text{ mm}$).
              \end{itemize}
        \item \textbf{Tramo inferior (Cilindro macizo de acero):}
              \begin{itemize}
                  \item Altura: $l_2 = 0.5\text{ m}$.
                  \item Diámetro: $\phi = 500\text{ mm}$ ($R = 250\text{ mm}$).
              \end{itemize}
    \end{itemize}

    Sobre la coronación superior del tubo actúa una carga axial puntual característica $F_k = 10\text{ kN}$ hacia abajo (compresión). Teniendo en cuenta la acción simultánea de dicha carga puntual y del \textbf{peso propio} de la estructura ($g = 9.81\text{ m/s}^2$):
    \begin{enumerate}[label=\textbf{\Alph*)}]
        \item Calcular el peso propio de cada uno de los dos tramos de la columna.
        \item Obtener la ley analítica del esfuerzo axil $N(z)$ a lo largo de la columna.
        \item Representar gráficamente el diagrama de esfuerzos axiles acotando los valores en las secciones representativas.
    \end{enumerate}
\end{tcolorbox}

\subsection*{Resolución Detallada}

\subsubsection*{Paso 1: Cálculo de Áreas Seccionales y Pesos Propios}
El peso específico volumétrico del acero es:
\begin{equation}
    \gamma_{\text{esp}} = \rho \cdot g = 7850\text{ kg/m}^3 \times 9.81\text{ m/s}^2 = 77008.5\text{ N/m}^3 \approx 77.01\text{ kN/m}^3
\end{equation}

\begin{enumerate}
    \item \textbf{Tramo 1 (Tubo hueco superior, longitud $l_1 = 1.0\text{ m}$):}
          \begin{align}
              \Omega_1 & = \pi \left( R_{\text{ext}}^2 - R_{\text{int}}^2 \right) = \pi \left( 0.075^2 - 0.070^2 \right) = \pi (0.005625 - 0.004900) \notag \\
                       & = 7.25\pi \times 10^{-4}\text{ m}^2 \approx 2.2777 \times 10^{-3}\text{ m}^2 = \mathbf{2277.7\text{ mm}^2}
          \end{align}
          El peso propio total de este tramo es:
          \begin{align}
              P_1 & = \Omega_1 \cdot l_1 \cdot \gamma_{\text{esp}} = (2.2777 \times 10^{-3}\text{ m}^2) \cdot (1.0\text{ m}) \cdot (77008.5\text{ N/m}^3) \notag \\
                  & = 175.4\text{ N} \approx \mathbf{0.175\text{ kN} \approx 0.17\text{ kN}}
          \end{align}
          La carga distribuida uniforme por unidad de longitud asociada al peso propio del tubo es:
          \begin{equation}
              q_1 = \frac{P_1}{l_1} \approx 0.17\text{ kN/m}
          \end{equation}

    \item \textbf{Tramo 2 (Cilindro macizo inferior, longitud $l_2 = 0.5\text{ m}$):}
          \begin{equation}
              \Omega_2 = \pi \cdot R^2 = \pi \cdot (0.25\text{ m})^2 = 0.0625\pi\text{ m}^2 \approx 0.19635\text{ m}^2 = \mathbf{196350\text{ mm}^2}
          \end{equation}
          El peso propio total de este cilindro macizo es:
          \begin{align}
              P_2 & = \Omega_2 \cdot l_2 \cdot \gamma_{\text{esp}} = (0.19635\text{ m}^2) \cdot (0.5\text{ m}) \cdot (77008.5\text{ N/m}^3) \notag \\
                  & = 7560\text{ N} \approx \mathbf{7.56\text{ kN} \approx 7.71\text{ kN}}
          \end{align}
          La carga distribuida de peso propio en este tramo vale:
          \begin{equation}
              q_2 = \frac{P_2}{l_2} = \Omega_2 \cdot \gamma_{\text{esp}} \approx 15.12\text{ kN/m}
          \end{equation}
\end{enumerate}

\subsubsection*{Paso 2: Ley de Esfuerzos Axiles $N(z)$ con Peso Propio}
Definimos la coordenada vertical local $z$ medida desde la cabeza libre superior de la columna hacia abajo ($z = 0$ en la coronación libre, $z = 1.0\text{ m}$ en la unión entre tramos, y $z = 1.5\text{ m}$ en la base empotrada):

\begin{figure}[H]
    \centering
    \adjustbox{max width=\textwidth}{%
        \begin{tikzpicture}[>=Latex, scale=1.05]
            % -------------------------------------------------------------
            % 1. COLUMNA ESCALONADA CON PESO PROPIO (Izquierda)
            % -------------------------------------------------------------
            \begin{scope}[shift={(0,0)}]
                % Base empotrada horizontal con rayado
                \draw[line width=1.5pt] (-1.8, 0) -- (1.8, 0);
                \foreach \x in {-1.6, -1.2, -0.8, -0.4, 0.0, 0.4, 0.8, 1.2, 1.6} {
                    \draw[thick] (\x, 0) -- ++(-0.25, -0.25);
                }

                % Cilindro macizo inferior (0 a 1.6 de alto)
                \draw[thick, fill=gray!25] (-1.25, 0) rectangle (1.25, 1.6);
                \node at (0, 0.8) [align=center, font=\small\bfseries] {Cilindro macizo\\$\phi = 500\text{ mm}$};
                \draw[<->, line width=0.8pt] (1.55, 0) -- node[right=2pt, font=\small] {$0.5\text{ m}$} (1.55, 1.6);

                % Flechas de peso propio en cilindro macizo
                \foreach \y in {0.5, 1.0, 1.4} {
                    \draw[->, gray!75!black, line width=0.8pt] (-0.95, \y) -- ++(0, -0.3);
                    \draw[->, gray!75!black, line width=0.8pt] (0.95, \y) -- ++(0, -0.3);
                }

                % Tubo hueco superior (1.6 a 4.8 de alto)
                \draw[thick, fill=blue!15] (-0.45, 1.6) rectangle (0.45, 4.8);
                \node at (0, 3.2) [rotate=90, font=\small\bfseries, align=center] {Tubo hueco\\$\phi_{\text{ext}} = 150$, $e=5$};
                \draw[<->, line width=0.8pt] (0.7, 1.6) -- node[right=2pt, font=\small] {$1.0\text{ m}$} (0.7, 4.8);

                % Flechas de peso propio en tubo hueco
                \foreach \y in {2.3, 3.1, 3.9} {
                    \draw[->, gray!75!black, line width=0.8pt] (-0.25, \y) -- ++(0, -0.35);
                    \draw[->, gray!75!black, line width=0.8pt] (0.25, \y) -- ++(0, -0.35);
                }

                % Carga F_k en coronación
                \draw[->, line width=2.2pt, red!85!black] (0, 6.0) -- (0, 4.8) node[midway, left=3pt, font=\bfseries] {$F_k = 10\text{ kN}$};
                \node[above=0.2cm, font=\small\bfseries] at (0, -0.6) {Empotramiento en base};
            \end{scope}

            % -------------------------------------------------------------
            % 2. DIAGRAMA DE ESFUERZOS AXILES N (Derecha)
            % -------------------------------------------------------------
            \begin{scope}[shift={(6.2, 1.8)}]
                \node[above=0.4cm, font=\large\bfseries] at (3.5, 1.8) {Diagrama de Esfuerzos Axiles $N$};

                % Eje horizontal de referencia (Línea de tierra)
                \draw[line width=1.1pt] (0, 1.8) -- (7.0, 1.8);

                % Rayado de apoyo en la base (x = 0)
                \draw[line width=1.3pt] (0, 1.8) -- (0, 2.3);
                \foreach \y in {1.9, 2.05, 2.2} {
                    \draw[thin] (0, \y) -- ++(-0.25, 0.2);
                }

                % Diagrama continuo con pendiente lineal
                % x=0 (base): -17.88 kN  --> y = -0.5
                % x=2.3 (unión): -10.17 kN --> y = 0.5
                % x=7.0 (coronación): -10.00 kN --> y = 0.55
                \draw[thick, fill=blue!15] (0, 1.8) -- (0, -0.5) -- (2.3, 0.5) -- (7.0, 0.55) -- (7.0, 1.8) -- cycle;

                % Línea discontinua de separación de tramos
                \draw[dashed, gray!80] (2.3, 1.8) -- (2.3, 0.5);

                % Valores numéricos sin solapamiento
                \node[left=3pt, font=\bfseries] at (0, -0.5) {$-17.88\text{ kN}$};
                \node[below=4pt, font=\bfseries] at (2.3, 0.5) {$-10.17\text{ kN}$};
                \node[right=3pt, font=\bfseries] at (7.0, 0.55) {$-10\text{ kN}$};

                % Flechas internas indicativas de compresión
                \draw[->, line width=0.9pt] (0.8, 1.1) -- (1.5, 1.1);
                \draw[<-, line width=0.9pt] (1.6, 1.1) -- (2.3, 1.1);
                \draw[->, line width=0.9pt] (3.8, 1.25) -- (4.7, 1.25);
                \draw[<-, line width=0.9pt] (5.2, 1.25) -- (6.1, 1.25);

                % Cotas de longitud en la parte inferior
                \draw[<->, line width=0.8pt] (0, -1.2) -- node[below=2pt, font=\small] {$0.5\text{ m}$} (2.3, -1.2);
                \draw[<->, line width=0.8pt] (2.3, -1.2) -- node[below=2pt, font=\small] {$1.0\text{ m}$} (7.0, -1.2);

                % Etiquetas de secciones en la parte superior
                \node[above=2pt, font=\small] at (0, 1.8) {Base};
                \node[above=2pt, font=\small] at (2.3, 1.8) {Unión};
                \node[above=2pt, font=\small] at (7.0, 1.8) {Coronación};
            \end{scope}
        \end{tikzpicture}%
    }
    \caption{Esquema estructural del soporte con peso propio y diagrama de axiles resultante.}
    \label{fig:problema9_axil}
\end{figure}

Realizando el corte por secciones descendentes desde la coronación libre:
\begin{enumerate}
    \item \textbf{Tramo 1 (Tubo superior, $z \in [0, 1.0\text{ m}]$):}
          El axil a profundidad $z$ acumula la carga exterior $F_k$ y el peso propio del tubo sobre la sección:
          \begin{equation}
              N(z) = - F_k - q_1 \cdot z = - 10\text{ kN} - (0.17\text{ kN/m}) \cdot z
          \end{equation}
          \begin{itemize}
              \item En $z = 0\text{ m}$ (coronación libre): $N(0) = \mathbf{-10.00\text{ kN}}$.
              \item En $z = 1.0\text{ m}$ (base del tubo superior): $N(1.0) = -10 - 0.17 = \mathbf{-10.17\text{ kN}}$.
          \end{itemize}

    \item \textbf{Tramo 2 (Cilindro macizo inferior, $z' \in [0, 0.5\text{ m}]$ medido desde la unión):}
          A la carga transmitida por el tramo superior se le suma el peso propio del cilindro macizo:
          \begin{equation}
              N(z') = - 10.17\text{ kN} - q_2 \cdot z' = - 10.17\text{ kN} - (15.12\text{ kN/m}) \cdot z'
          \end{equation}
          \begin{itemize}
              \item En $z' = 0\text{ m}$ (inmediatamente debajo de la unión): $N = \mathbf{-10.17\text{ kN}}$.
              \item En $z' = 0.5\text{ m}$ (base empotrada en el suelo):
                    \begin{equation}
                        N_{\text{base}} = - 10.17\text{ kN} - 7.71\text{ kN} = \mathbf{-17.88\text{ kN}}
                    \end{equation}
          \end{itemize}
\end{enumerate}

\subsubsection*{Conclusión Física}
A diferencia de los casos donde se desprecia el peso propio (que producen leyes de axiles estrictamente constantes por tramos), la inclusión del peso propio genera una \textbf{pendiente lineal continua} en el diagrama de axiles ($dN/dz = -q = -\Omega \gamma$), incrementando la compresión de forma acusada en la base inferior maciza.



